% 原书第142—148页：7.2导言及7.2.1。
\FloatBarrier
\section{结构识别方法}

前面介绍了几种能够使因果结构具有全部或部分可识别性的假设。本节说明如何利用这些假设，根据有限数据构造底层图的估计量，图7.1展示了两种思路。这里对各类方法作概述，重点介绍其思想。目前可用的方法很多，哪些方法在实际中最有价值，我们认为还有待今后的研究检验。不过，我们仍会尽量指出其中一些方法的潜在问题和最关键的假设。虽然部分文献研究了这些方法的一致性，但这里略去大多数这类结果，只介绍思路。算法实现中的细节也不作讨论，有兴趣的读者可参阅所列文献。Kalisch et al.［2012］维护的 R 软件包 \texttt{pcalg}［R Core Team, 2016］，不仅包含 PC 算法的代码，名称来自其提出者 Peter Spirtes 和 Clark Glymour，见第7.2.1节，也实现了本节介绍的许多其他方法。

在详细介绍现有方法之前，先补充两点说明。
\begin{enumerate}[label=(\arabic*)]
 \item 虽然已有多项模拟研究，但损失函数问题受到的关注仍然很少。给定真实的底层因果结构，应如何衡量估计因果图的“好坏”？实际中，人们常采用\textbf{结构 Hamming 距离（structural Hamming distance）}的各种变体［Acid and de Campos, 2003；Tsamardinos et al., 2006］，统计指定错误的边数。另一种选择是依据图预测干预分布的能力来评价它［Peters and Bühlmann, 2015］。
 \item 本节介绍的某些方法，假定结构赋值 (6.1)，特别是相应的函数 $f_j$，具有简单形式。这些方法往往不仅估计因果结构，还估计对应的赋值，而后者通常也可用于计算残差。原则上，在这一模型下，就可以检验“噪声变量相互独立”这一很强的假设，见定义3.1。例如，可以应用相互独立性检验［如 Pfister et al., 2017］；这种过程在统计上涉及的细节，见第4.2.1节。
\end{enumerate}

% 原图7.1：逐节点、逐箭头核对原书第144页后以中文TikZ重排。
\usetikzlibrary{shapes.symbols}
\begin{figure}[htbp]
\centering
\begin{tikzpicture}[causal,x=1cm,y=1cm,font=\small,
 chsevenbox/.style={draw,rounded corners=2mm,align=left,inner sep=7pt},
 chsevenflow/.style={->,semithick}]
 \node[cloud,draw,cloud puffs=9,cloud puff arc=110,aspect=1.8,
       fill=gray!12,align=center,inner sep=3pt,minimum width=2.5cm,
       minimum height=1.2cm] (ch07f1sample) at (1.2,3.7)
       {来自 $P_{X_1,\ldots,X_4}$ 的\\独立同分布样本};
 \node[chsevenbox] (ch07f1indep) at (6.25,5.85)
       {$X_2\indep(X_1,X_3)$\\[3pt]
        $X_1\indep X_4\mid X_3$\\[3pt]
        $X_1\indep X_2\mid X_3$\\[3pt]
        $X_2\indep X_3\mid X_1$\\[2pt]
        \hspace{1cm}$\cdots$};
 \node[chsevenbox] (ch07f1scm) at (6.25,1.35)
       {$X_1:=f_1(N_1)$\\[3pt]
        $X_2:=f_2(N_2)$\\[3pt]
        $X_3:=f_3(X_1,N_3)$\\[3pt]
        $X_4:=f_4(X_2,X_3,N_4)$\\[3pt]
        $N_1,\ldots,N_4$ 相互独立};
 \draw[rounded corners=2mm,fill=gray!12]
       (10.0,2.35) rectangle (13.8,5.05);
 \node (ch07f1glabel) at (10.65,4.55) {$\mathcal G$};
 \node[vtx,fill=white] (ch07f1x1) at (11.8,4.55) {$X_1$};
 \node[vtx,fill=white] (ch07f1x2) at (10.7,3.65) {$X_2$};
 \node[vtx,fill=white] (ch07f1x3) at (13.0,3.65) {$X_3$};
 \node[vtx,fill=white] (ch07f1x4) at (11.8,2.85) {$X_4$};
 \draw[ed] (ch07f1x1)--(ch07f1x3);
 \draw[ed] (ch07f1x2)--(ch07f1x4);
 \draw[ed] (ch07f1x3)--(ch07f1x4);
 \draw[chsevenflow] (ch07f1sample.north)
       .. controls (1.2,5.9) and (3.1,5.85) .. (ch07f1indep.west);
 \node[align=center] (ch07f1testlabel) at (3.1,6.3) {独立性检验};
 \draw[chsevenflow] (ch07f1indep.east)--(13.2,5.85);
 \node (ch07f1gprime) at (13.65,5.85) {$\mathcal G'$};
 \draw[chsevenflow] (10.9,5.85)
       .. controls (11.4,5.85) and (11.8,5.5) .. (11.9,5.15);
 \node[align=center] (ch07f1assumptions) at (10.55,6.35)
       {忠实性与\\马尔可夫条件};
 \draw[chsevenflow] (ch07f1sample.south)
       .. controls (1.2,1.35) and (3.1,1.35) .. (ch07f1scm.west);
 \node[align=center] (ch07f1fitlabel) at (2.9,0.9) {拟合受限\\模型类};
 \draw[chsevenflow] (ch07f1scm.east)
       .. controls (10.2,1.35) and (11.2,1.45) .. (11.9,2.23);
 \node (ch07f1readlabel) at (10.15,0.9) {读取图结构};
\end{tikzpicture}
\caption{识别因果结构的两种途径。基于独立性的方法（上方）检验数据中的条件独立关系，再通过马尔可夫条件和忠实性将这些性质与图结构联系起来。图往往不能唯一识别，因此，方法可能输出不同的图 $\mathcal G,\mathcal G'$。另一种途径是限制模型类，直接拟合结构因果模型（下方）。}
\label{fig:ch07-identification}
\end{figure}


\subsection{基于独立性的方法}

归纳因果算法（inductive causation，IC）、SGS 算法，名称来自提出者 Spirtes、Glymour 和 Scheines，以及 PC 算法等基于独立性的方法，都假定分布忠实于底层有向无环图。这使得马尔可夫等价类，即相应的 CPDAG，具有可识别性，见第7.1.1节。图中的 d-分离与 $P_X$ 中的条件独立关系之间存在一一对应。因此，对于任意 d-分离陈述，都可以通过相应的条件独立性检验来作出判断。先假定一个\textbf{预言机（oracle）}能够正确回答条件独立性问题；有限样本下的问题将在下面“条件独立性检验”一段讨论。

\paragraph{估计骨架}
大多数基于独立性的方法先估计骨架，也就是无向边，再尽可能为边确定方向。对于骨架搜索，下面的引理很有用［见 Verma and Pearl, 1991，引理1］。

\begin{lemma}
下面两条结论成立。
\begin{enumerate}[label=(\roman*)]
 \item 有向无环图 $(\mathbf X,\mathcal E)$ 中的两个节点 $X,Y$ 相邻，当且仅当任何子集 $S\subseteq V\setminus\{X,Y\}$ 都不能将它们 d-分离。
 \item 如果有向无环图 $(\mathbf X,\mathcal E)$ 中的两个节点 $X,Y$ 不相邻，那么 $\PA_X$ 或 $\PA_Y$ 至少有一个能将它们 d-分离。
\end{enumerate}
\end{lemma}

由引理7.8（i）可知，如果无论以其他哪些变量为条件，两个变量总是相依，那么它们一定相邻。IC 算法［Pearl, 2009］和 SGS 算法［Spirtes et al., 2000］利用了这一结果。对于每对节点 $(X,Y)$，这些方法遍历所有既不包含 $X$ 也不包含 $Y$ 的变量子集 $A\subseteq\mathbf X\setminus\{X,Y\}$，检查给定 $A$ 后 $X$ 与 $Y$ 是否 d-分离。完成所有检验后，$X,Y$ 相邻，当且仅当没有找到能将它们 d-分离的集合 $A$。

遍历所有可能的子集 $A$ 似乎并不理想，图稀疏时尤其如此。\textbf{PC算法}［Spirtes et al., 2000］从完全连接的无向图出发，从 $\#A=0$ 开始，逐步增加条件集合 $A$ 的大小。在第 $k$ 次迭代中，它考虑大小为 $\#A=k$ 的集合，并使用一个巧妙技巧：要检验 $X$ 与 $Y$ 是否能被 d-分离，只需遍历那些完全包含在 $X$ 的邻节点集合或者 $Y$ 的邻节点集合中的 $A$。这一思路基于引理7.8（ii），显然能够缩短计算时间，尤其适用于稀疏图。

\paragraph{确定边的方向}
引理6.25表明，应当能够确定图中非道德结构（immoralities），也就是 v 结构的方向。如果所得骨架中的两个节点之间没有直接连边，那么就存在一个集合将它们 d-分离。假设骨架包含结构 $X-Z-Y$，而 $X,Y$ 之间没有直接边；再设 $A$ 是一个使 $X,Y$ d-分离的集合。此时，$X-Z-Y$ 是非道德结构，因而能够定向为 $X\to Z\leftarrow Y$，当且仅当 $Z\notin A$。确定这些结构的方向后，出于避免有向环等要求，还可能确定某些其他边的方向。一组这样的定向规则已被证明具有完备性，称为\textbf{Meek定向规则（Meek's orientation rules）}［Meek, 1995］。

\paragraph{可满足性方法}
除了上述图形方法，还可以把因果学习表述为\textbf{可满足性问题（satisfiability，SAT）}［Triantafillou et al., 2010］。首先，用布尔变量表示图关系，例如定义 $A$ 为“存在从 $X$ 到 $Y$ 的直接边”。困难之处在于，将独立性陈述转化为 d-分离陈述，再将其翻译成包含布尔变量及“与”“或”运算符的公式；这里仍假设独立性陈述由独立性预言机给出。SAT 问题询问的是：能否为每个布尔变量赋值为“真”或“假”，使整个公式为真？

SAT 求解器不仅检查是否存在这样的赋值，还可以告诉我们：在所有能使整体公式为真的赋值中，某些变量是否始终取同一个值。例如，不同图结构可能都满足相同的 d-分离陈述，并对应不同的布尔赋值；但如果在所有这样的赋值中，上述布尔变量 $A$ 都取“真”，就可以推断，底层图中的 $X$ 必须是 $Y$ 的父节点。虽然布尔 SAT 问题已知是\textbf{NP完全}的［Cook, 1971；Levin, 1973］，即既属于非确定性多项式时间（NP）问题，又是 NP 难的，但仍有启发式算法可以求解包含数百万个变量的大型问题实例。

因果学习中的 SAT 方法允许查询祖先关系之类的具体陈述，而不必估计完整图。它们能够纳入不同类型的先验知识；如果认为某些统计发现相互矛盾，还可以为独立性约束赋予不同权重。这类方法已被推广到有向环、潜变量以及相互重叠的数据集等情形［Hyttinen et al., 2013；Triantafillou and Tsamardinos, 2015］。

\paragraph{条件独立性检验}
在前面的三段讨论中，我们假设存在一个独立性预言机，能够告诉我们，分布中是否存在某个特定的（条件）独立关系。然而，在实际中，必须根据有限数据推断这一陈述。这带来两项主要挑战。
\begin{enumerate}[label=(\arabic*)]
 \item 所有基于条件独立性检验的因果发现方法，都会同时依据相依性与独立性作出结论。但在实际中，最常用的是统计显著性检验，而它本身是不对称的。因此，人们通常会搁置显著性水平原本的含义，把它当作一个调节参数。此外，由于样本有限，检验结果甚至可能相互矛盾：不存在任何图结构能够精确编码所推断出的条件独立关系集合。
 \item 虽然近来出现了一些基于核的检验［Fukumizu et al., 2008；Tillman et al., 2009；Zhang et al., 2011］，但在数据量有限时，非参数条件独立性检验很难实施。因此，人们常将讨论限制在可能相依关系的某个子类内。下面简要介绍其中几种情形。
\end{enumerate}

如果假定变量服从高斯分布，可以检验偏相关是否为零，见附录A.1与A.2。忠实性使底层有向无环图的马尔可夫等价类具有可识别性，见引理7.2。事实上，在高斯情形下，使用零偏相关检验的 PC 算法，能够给出正确 CPDAG 的一致估计［Kalisch and Bühlmann, 2007］。如果再假设一个称为\textbf{强忠实性（strong faithfulness）}的条件［Zhang and Spirtes, 2003；Uhler et al., 2013］，甚至可以得到\textbf{均匀一致性（uniform consistency）}［Kalisch and Bühlmann, 2007］；另见 Robins et al.［2003］中的讨论。

无论在理论上还是实际中，非参数条件独立性检验都很困难。对于非高斯分布，零偏相关既不是条件独立的必要条件，也不是充分条件，下面的例子说明了这一点。

\begin{example}[条件独立与偏相关]
\begin{enumerate}[label=(\roman*)]
 \item 如果分布 $P_{X,Y,Z}$ 由下列结构因果模型导出：
 \[
  Z:=N_Z,\qquad X:=Z^2+N_X,\qquad Y:=Z^2+N_Y,
 \]
 其中 $N_X,N_Y,N_Z\overset{\mathrm{i.i.d.}}{\sim}\mathcal N(0,1)$，那么它满足
 \[
  X\indep Y\mid Z,\qquad \rho_{X,Y\mid Z}\neq0.
 \]
 偏相关系数 $\rho_{X,Y\mid Z}$ 等于 $X-\alpha Z$ 与 $Y-\beta Z$ 的相关系数，其中 $\alpha,\beta$ 分别为 $X$ 关于 $Z$ 和 $Y$ 关于 $Z$ 回归所得的回归系数。在本例中，$\alpha=\beta=0$，因为 $X$ 和 $Y$ 都与 $Z$ 不相关。
 \item 下列结构因果模型导出的分布 $P_{X,Y,Z}$，
 \[
  Z:=N_Z,\qquad X:=Z+N_X,\qquad Y:=Z+N_Y,
 \]
 其中 $(N_X,N_Y)\indep N_Z$，且 $N_X,N_Y$ 不相关但不独立，满足
 \[
  X\nindep Y\mid Z,\qquad \rho_{X,Y\mid Z}=0,
 \]
 因为此时 $\rho_{X,Y\mid Z}$ 就是 $N_X$ 与 $N_Y$ 的相关系数。
\end{enumerate}
因此，零偏相关既不蕴含条件独立，也不由条件独立所蕴含。
\end{example}

下面这个检验给定 $Z$ 后 $X$ 与 $Y$ 是否条件独立的过程，是偏相关的一种自然非线性推广［例如 Ramsey, 2014］：
\begin{enumerate}[label=(\arabic*)]
 \item 对 $X$ 关于 $Z$ 作非线性回归，检验残差是否与 $Y$ 独立；
 \item 对 $Y$ 关于 $Z$ 作非线性回归，检验残差是否与 $X$ 独立；
 \item 如果上述两种独立关系中至少有一种成立，就得出 $X\indep Y\mid Z$。
\end{enumerate}
在加性噪声模型情形下，这似乎是正确的检验，见第7.1.2节。例如，对于三个变量，有以下结论。

\begin{proposition}
考虑由加性噪声模型，即定义7.3，诱导的分布 $P_{X,Y,Z}$，其中所有变量都具有严格为正的密度。如果给定 $Z$ 后 $X$ 与 $Y$ d-分离，那么上述过程会输出相应的条件独立关系；具体地，$X-\EE[X\mid Z]$ 与 $Y$ 独立，或者 $Y-\EE[Y\mid Z]$ 与 $X$ 独立，二者至少成立其一。
\end{proposition}
\begin{proof}
假设 $X:=h(Z)+N_X$、$Y:=f(Z)+N_Y$，且 $Z,N_X,N_Y$ 相互独立。那么，$X-\EE[X\mid Z]=N_X$ 与 $Y$ 独立。\chseventnote{原书将残差写成 $N_X$，这隐含了噪声均值为零。一般应为 $N_X-\EE[N_X]$，但减去常数不影响其与 $Y$ 的独立性，因此本命题的独立性结论不受影响。这些含期望的表达式还默认相应期望存在。}
对于其他可能的结构，例如 $X\to Z\to Y$ 或 $X\leftarrow Z\leftarrow Y$，可以类似地得到结论。
\end{proof}

这个命题说明，从总体层面看，上述检验适用于三个变量的加性噪声模型。然而，考虑四个变量 $X,Y,Z,V$ 时，就已经可能出现问题。\chseventnote{原书先列四个变量 $X,Y,Z,V$，接着的图结构却使用 $W$。这里的 $V$ 与 $W$ 应理解为同一个第四变量，正文保留原书记号。}
显然，图 $X\leftarrow Z\to W\to Y$ 与 $X\to Z\to W\to Y$ 马尔可夫等价。不过，虽然该检验对第一个图会输出 $X\indep Y\mid Z$，对第二个图却没有这种保证。因此，为构造可实施的条件独立性检验而对随机变量之间相依模型作出的上述限制，会导致同一个马尔可夫等价类内的图受到不对称处理。许多其他类型的条件独立性检验，也可能存在同样的效应。这种不对称性不一定是缺点，因为正如前面看到的，受限函数类可能使同一个马尔可夫等价类内的图也具有可识别性，见第7.1节。不过，这一问题无疑需要认真考虑。
